% ---- pg-552 （书页 540） ----

集中考察中间层的聚类, 我们发现公用事业公司倾向于按地理位置分组. 例如, 一个中间聚类包含企业 1(亚利桑那公共服务公司), 18(南方公司——业务范围在佐治亚和阿拉巴马), 19(得克萨斯公用事业公司)及 14(俄克拉何马天然气与电力公司). 也有某些例外. 聚类 $(7,12,21,15,2)$ 包含沿东海岸和远西的企业. 另一方面, 所有这些企业都位于沿海地区. 请注意: 纽约的联合爱迪生公司和圣地亚哥天然气与电力公司直到最后的大汇合之前均自成一个聚类.

也许, 位置相似(或位置的类型相似)的公用事业公司构成一个聚类是不值得惊奇的. 可以预料: 处在同一地区的公司其发电厂基本上使用同一类型的燃料, 而且面对同一个市场. 因此, 发电的类型、成本、增长等在这些企业之间是相对一致的. 这些内容都被表观地反映在分层聚类方法之中. \hfill$\blacksquare$

\noindent 对平均连接聚类, 距离(相似性)的改变会影响最终聚类结构, 即使这种改变是保序的.

\mvsub{沃德(Ward)的分层聚类法}

沃德(见文献[32])考虑了建立在联合两个聚类时使“信息损失”最小这一想法基础上的分层聚类方法.
这种方法通常在实施过程中会有信息损失, 用某种误差平方和(ESS)准则来衡量时表现为 ESS 有所增大.
首先, 对给定的聚类 $k$, 令 $\mathrm{ESS}_k$ 为此聚类中每一个项目与类均值(中心点)的偏差平方和.
若当前有 $K$ 个聚类, 则定义 $\mathrm{ESS}=\mathrm{ESS}_1+\mathrm{ESS}_2+\cdots+\mathrm{ESS}_K$. 在分析过程的每一步, 考虑可能有的各对聚类的合并.
若某对聚类的合并使得 ESS 增加最少(信息损失最小), 则将这两个聚类合并.
最初, 每个聚类仅由一个项目组成, 因而如果有 $N$ 个项目, 则有 $\mathrm{ESS}_k=0$, $k=1,2,\cdots,N$, 故 $\mathrm{ESS}=0$.
在另一极端, 当所有聚类合并成一个由 $N$ 个项目组成的组, 则 ESS 的值为
\[
\mathrm{ESS}=\sum_{j=1}^{N}(\mathbf{x}_j-\bar{\mathbf{x}}) '(\mathbf{x}_j-\bar{\mathbf{x}})
\]
其中 $\mathbf{x}_j$ 为与第 $j$ 个项目相联系的多元测量值, 而 $\bar{\mathbf{x}}$ 则为总均值.

沃德方法的结果可以用树图表示, 树图的纵坐标代表发生合并时的 ESS 值.

沃德方法基于这样一个概念, 即预期多元观测值的聚类大体呈椭球形分布. 它是非分层聚类方法的先驱.
非分层聚类法通过优化某种准则, 实现将数据分割为给定数目的椭圆状聚类的目的.
我们将在下一节中讨论非分层聚类法. 关于聚类分析中优化方法的进一步讨论可以在文献[8]中找到.

\noindent\textbf{例 12.10}\quad (纯麦芽苏格兰威士忌的聚类)

世界上所有的纯麦芽苏格兰威士忌事实上均产于苏格兰. 在一项研究(见文献[22])中, 曾用 68 个二变量来测量苏格兰威士忌的各种特性.
这些变量大致可分为颜色、气味、酒质、口感以及最后加工几类. 例如, 有 14 种颜色特性(描述), 其中包括白葡萄调色度、黄色、极浅色、浅色、红褐色、金琥珀色, 红色等.
拉普安特(LaPointe)和勒让德(Legendre)将来自不同酿酒厂的 109 种纯麦芽威士忌(均来自不同酒厂)进行了聚类.
研究者的兴趣在于确定纯麦芽苏格兰威士忌的主要类型, 它们的主要特性, 以及最佳代表.
此外, 他们想知道由分层聚类法所得到的聚类是否与不同的地理区域相对应, 因为已经知道威士忌受土壤、温度和水质这些条件的影响.

用二变量表示特性的存在或缺失, 在此基础上产生加权的相似系数 $(S_{ik})$, 而“距离”定义为 $d_{ik}=1-S_{ik}$.
用沃德方法对 109 种纯(单一)麦芽苏格兰威士忌分类, 得到的树图示于\hyperref[fig:12.11]{图 12.11}. (将平均连接法用于相似性矩阵, 几乎得出完全相同的分类.)
